\documentclass[12pt,a4paper]{article}

\usepackage[T2A]{fontenc}
\usepackage[utf8]{inputenc}
\usepackage[russian]{babel}
\usepackage{textcomp}
\usepackage{amsmath,amssymb}
\usepackage{array}
\usepackage{enumitem}
\usepackage{xcolor}
\usepackage{listings}
\usepackage{geometry}
\usepackage{hyperref}

\geometry{left=2.5cm,right=2.0cm,top=2.0cm,bottom=2.0cm}
\emergencystretch=2em

\lstdefinelanguage{JavaScript}{
    keywords={
        const, let, var, function, return, if, else, for, while, do, break,
        continue, switch, case, default, new, class, extends, import, export,
        from, async, await, try, catch, finally, throw, document, window
    },
    sensitive=true,
    comment=[l]{//},
    morecomment=[s]{/*}{*/},
    morestring=[b]',
    morestring=[b]",
    morestring=[b]`
}

\lstdefinelanguage{CSS}{
    keywords={
        color, background, font-family, max-width, margin, padding, display,
        flex, gap, flex-direction, column, border-radius
    },
    sensitive=true,
    comment=[s]{/*}{*/},
    morestring=[b]',
    morestring=[b]"
}

\lstdefinelanguage{HTML}{
    keywords={
        html, head, body, title, meta, link, script, h1, h2, p, a, div, span,
        header, nav, main, section, article, aside, footer, table, tr, th, td,
        form, label, input, button, ol, ul, li, strong, em
    },
    sensitive=false,
    comment=[s]{<!--}{-->},
    morestring=[b]",
    morestring=[b]'
}

\lstdefinelanguage{Python}{
    keywords={
        import, from, as, def, return, if, elif, else, for, while, break,
        continue, class, try, except, finally, with, True, False, None, print
    },
    sensitive=true,
    comment=[l]{\#},
    morestring=[b]',
    morestring=[b]"
}

\lstdefinestyle{code}{
    basicstyle=\ttfamily\small,
    breaklines=true,
    frame=single,
    columns=fullflexible,
    keepspaces=true,
    showstringspaces=false,
    inputencoding=utf8,
    extendedchars=true,
    literate=
        {А}{{\CYRA}}1 {Б}{{\CYRB}}1 {В}{{\CYRV}}1 {Г}{{\CYRG}}1
        {Д}{{\CYRD}}1 {Е}{{\CYRE}}1 {Ё}{{\CYRYO}}1 {Ж}{{\CYRZH}}1
        {З}{{\CYRZ}}1 {И}{{\CYRI}}1 {Й}{{\CYRISHRT}}1 {К}{{\CYRK}}1
        {Л}{{\CYRL}}1 {М}{{\CYRM}}1 {Н}{{\CYRN}}1 {О}{{\CYRO}}1
        {П}{{\CYRP}}1 {Р}{{\CYRR}}1 {С}{{\CYRS}}1 {Т}{{\CYRT}}1
        {У}{{\CYRU}}1 {Ф}{{\CYRF}}1 {Х}{{\CYRH}}1 {Ц}{{\CYRC}}1
        {Ч}{{\CYRCH}}1 {Ш}{{\CYRSH}}1 {Щ}{{\CYRSHCH}}1 {Ъ}{{\CYRHRDSN}}1
        {Ы}{{\CYRERY}}1 {Ь}{{\CYRSFTSN}}1 {Э}{{\CYREREV}}1
        {Ю}{{\CYRYU}}1 {Я}{{\CYRYA}}1
        {а}{{\cyra}}1 {б}{{\cyrb}}1 {в}{{\cyrv}}1 {г}{{\cyrg}}1
        {д}{{\cyrd}}1 {е}{{\cyre}}1 {ё}{{\cyryo}}1 {ж}{{\cyrzh}}1
        {з}{{\cyrz}}1 {и}{{\cyri}}1 {й}{{\cyrishrt}}1 {к}{{\cyrk}}1
        {л}{{\cyrl}}1 {м}{{\cyrm}}1 {н}{{\cyrn}}1 {о}{{\cyro}}1
        {п}{{\cyrp}}1 {р}{{\cyrr}}1 {с}{{\cyrs}}1 {т}{{\cyrt}}1
        {у}{{\cyru}}1 {ф}{{\cyrf}}1 {х}{{\cyrh}}1 {ц}{{\cyrc}}1
        {ч}{{\cyrch}}1 {ш}{{\cyrsh}}1 {щ}{{\cyrshch}}1 {ъ}{{\cyrhrdsn}}1
        {ы}{{\cyrery}}1 {ь}{{\cyrsftsn}}1 {э}{{\cyrerev}}1
        {ю}{{\cyryu}}1 {я}{{\cyrya}}1
        {—}{{---}}1
        {–}{{--}}1
        {«}{{``}}1
        {»}{{''}}1
        {№}{{\textnumero}}1
        {…}{{\ldots}}1,
    keywordstyle=\color{blue},
    commentstyle=\color{gray},
    stringstyle=\color{red!60!black}
}

\lstset{style=code}

\title{Билет №29 \\ \small Удаленный доступ к ресурсам сети. Организация электронной почты, телеконференций. Протоколы передачи файлов FTP и HTTP, язык разметки гипертекста HTML, разработка WEB-страниц, WWW-серверы. (ИТМО)}
\author{Сериков Владислав Александрович}
\date{6 мая 2026}

\begin{document}

\maketitle
\tableofcontents
\newpage

\section{Базовые понятия компьютерных сетей}

\subsection{Ресурсы сети}

\textbf{Сетевой ресурс} --- это аппаратный, программный или информационный объект, доступный пользователям или программам через компьютерную сеть.

К сетевым ресурсам относятся:
\begin{itemize}
    \item файлы и каталоги на удалённых серверах;
    \item базы данных;
    \item почтовые ящики;
    \item веб-страницы и веб-приложения;
    \item вычислительные мощности удалённых машин;
    \item принтеры, сетевые хранилища, камеры и другие устройства;
    \item сервисы прикладного уровня: DNS, FTP, HTTP, SMTP, IMAP, SSH и др.
\end{itemize}

Доступ к ресурсам сети обычно осуществляется по модели ``клиент--сервер''.

\subsection{Модель клиент--сервер}

\textbf{Клиент} --- программа, инициирующая запрос к сетевому сервису.

\textbf{Сервер} --- программа или узел сети, принимающий запросы клиентов и предоставляющий им некоторый сервис.

Примеры:
\begin{center}
\begin{tabular}{|c|c|c|}
\hline
Сервис & Клиент & Сервер \\
\hline
Web & браузер & HTTP-сервер, например Apache или Nginx \\
\hline
Электронная почта & почтовый клиент & SMTP/IMAP/POP3-сервер \\
\hline
Передача файлов & FTP-клиент & FTP-сервер \\
\hline
Удалённое управление & SSH-клиент & SSH-сервер \\
\hline
\end{tabular}
\end{center}

\subsection{Адресация в сети}

Для доступа к ресурсу необходимо указать:
\begin{enumerate}
    \item адрес узла;
    \item протокол доступа;
    \item путь или идентификатор ресурса;
    \item при необходимости --- порт и параметры.
\end{enumerate}

В Internet для идентификации ресурсов используется URI.

\subsection{URI, URL и URN}

\textbf{URI} --- Uniform Resource Identifier, единообразный идентификатор ресурса.

\textbf{URL} --- Uniform Resource Locator, частный случай URI, указывающий не только имя ресурса, но и способ доступа к нему.

Общий вид URL:
\[
\text{scheme://host:port/path?query\#fragment}
\]

Пример:
\begin{lstlisting}
https://example.org:443/docs/index.html?id=10#section1
\end{lstlisting}

Здесь:
\begin{itemize}
    \item \texttt{https} --- схема, то есть протокол;
    \item \texttt{example.org} --- доменное имя узла;
    \item \texttt{443} --- порт;
    \item \texttt{/docs/index.html} --- путь к ресурсу;
    \item \texttt{id=10} --- строка запроса;
    \item \texttt{section1} --- фрагмент документа.
\end{itemize}

\subsection{DNS}

\textbf{DNS} --- Domain Name System, распределённая система доменных имён, преобразующая человекочитаемые имена в IP-адреса.

Например:
\[
\texttt{www.example.org} \rightarrow 93.184.216.34
\]

Упрощённый алгоритм разрешения доменного имени:
\begin{enumerate}
    \item Пользователь вводит адрес сайта в браузере.
    \item Браузер проверяет локальный кэш DNS.
    \item Если ответа нет, запрос передаётся системному DNS-резолверу.
    \item Резолвер обращается к корневым DNS-серверам.
    \item Затем --- к серверам домена верхнего уровня, например \texttt{.org}.
    \item Затем --- к авторитетному DNS-серверу домена.
    \item Полученный IP-адрес возвращается клиенту.
    \item Клиент устанавливает соединение с найденным узлом.
\end{enumerate}

Схематически:
\[
\text{браузер} \rightarrow \text{DNS-резолвер} \rightarrow \text{root} \rightarrow \text{TLD} \rightarrow \text{authoritative DNS}
\]

\section{Удалённый доступ к ресурсам сети}

\subsection{Понятие удалённого доступа}

\textbf{Удалённый доступ} --- это возможность взаимодействовать с ресурсами компьютера или сети, физически находясь в другом месте.

Удалённый доступ может быть:
\begin{itemize}
    \item интерактивным --- пользователь работает с удалённой системой почти как с локальной;
    \item файловым --- пользователь передаёт, получает или изменяет файлы;
    \item прикладным --- доступ к веб-приложению, почте, базе данных;
    \item административным --- настройка и сопровождение серверов.
\end{itemize}

\subsection{Основные способы удалённого доступа}

\begin{center}
\begin{tabular}{|p{3cm}|p{4cm}|p{6cm}|}
\hline
Способ & Протоколы/технологии & Назначение \\
\hline
Терминальный доступ & Telnet, SSH & Работа в командной строке удалённого компьютера \\
\hline
Удалённый рабочий стол & RDP, VNC & Графическое управление удалённой машиной \\
\hline
Файловый доступ & FTP, SFTP, SMB, NFS, WebDAV & Передача и совместное использование файлов \\
\hline
Web-доступ & HTTP, HTTPS & Доступ к сайтам и веб-приложениям \\
\hline
Почтовый доступ & SMTP, POP3, IMAP & Отправка и получение электронной почты \\
\hline
VPN & IPsec, OpenVPN, WireGuard & Создание защищённого канала в удалённую сеть \\
\hline
\end{tabular}
\end{center}

\subsection{Telnet и SSH}

\textbf{Telnet} --- старый протокол удалённого терминального доступа. Основной недостаток Telnet состоит в том, что данные, включая пароли, передаются в открытом виде.

\textbf{SSH} --- Secure Shell, защищённый протокол удалённого доступа. Он обеспечивает:
\begin{itemize}
    \item аутентификацию сервера;
    \item аутентификацию пользователя;
    \item шифрование трафика;
    \item контроль целостности данных;
    \item возможность туннелирования других протоколов.
\end{itemize}

Пример подключения:
\begin{lstlisting}
ssh student@example.org
\end{lstlisting}

Пример копирования файла через SSH:
\begin{lstlisting}
scp report.pdf student@example.org:/home/student/
\end{lstlisting}

\subsection{VPN}

\textbf{VPN} --- Virtual Private Network, виртуальная частная сеть. VPN создаёт защищённый логический канал поверх потенциально недоверенной сети.

Типичные задачи VPN:
\begin{itemize}
    \item доступ сотрудников к внутренней сети организации;
    \item защита трафика в публичных сетях;
    \item объединение филиалов;
    \item ограничение доступа к внутренним сервисам только через защищённый канал.
\end{itemize}

\subsection{Общие этапы удалённого доступа}

Типичный алгоритм удалённого доступа к ресурсу:
\begin{enumerate}
    \item Пользователь указывает адрес ресурса.
    \item Клиент определяет протокол доступа.
    \item При необходимости выполняется DNS-разрешение имени.
    \item Устанавливается сетевое соединение.
    \item Выполняется аутентификация сервера.
    \item Выполняется аутентификация пользователя, если она требуется.
    \item Клиент отправляет запрос.
    \item Сервер обрабатывает запрос.
    \item Сервер возвращает результат.
    \item Соединение закрывается или сохраняется для последующих запросов.
\end{enumerate}

\section{Электронная почта}

\subsection{Общее понятие}

\textbf{Электронная почта} --- система обмена электронными сообщениями между пользователями компьютерных сетей.

Адрес электронной почты имеет вид:
\[
\texttt{local-part@domain}
\]

Пример:
\begin{lstlisting}
student@example.org
\end{lstlisting}

Здесь:
\begin{itemize}
    \item \texttt{student} --- локальная часть адреса, обычно имя почтового ящика;
    \item \texttt{example.org} --- домен почтового сервиса.
\end{itemize}

\subsection{Основные компоненты почтовой системы}

\begin{enumerate}
    \item \textbf{MUA} --- Mail User Agent, пользовательский почтовый агент. Это почтовый клиент: Thunderbird, Outlook, Apple Mail, веб-интерфейс почты.
    \item \textbf{MSA} --- Mail Submission Agent, агент отправки почты от клиента к серверу.
    \item \textbf{MTA} --- Mail Transfer Agent, агент пересылки почты между серверами. Примеры: Postfix, Exim, Sendmail.
    \item \textbf{MDA} --- Mail Delivery Agent, агент доставки письма в почтовый ящик получателя.
    \item \textbf{Почтовое хранилище} --- место, где хранятся письма пользователя.
\end{enumerate}

\subsection{Протоколы электронной почты}

\begin{center}
\begin{tabular}{|c|c|c|}
\hline
Протокол & Назначение & Типичные порты \\
\hline
SMTP & Отправка и пересылка почты & 25, 465, 587 \\
\hline
POP3 & Получение почты с сервера & 110, 995 \\
\hline
IMAP & Синхронизированный доступ к почтовому ящику & 143, 993 \\
\hline
\end{tabular}
\end{center}

\subsection{SMTP}

\textbf{SMTP} --- Simple Mail Transfer Protocol, протокол передачи электронной почты.

SMTP применяется:
\begin{itemize}
    \item для отправки письма почтовым клиентом на сервер;
    \item для передачи письма между почтовыми серверами.
\end{itemize}

Упрощённый диалог SMTP:
\begin{lstlisting}
S: 220 mail.example.org ESMTP
C: EHLO client.example.com
S: 250-mail.example.org
S: 250 OK
C: MAIL FROM:<alice@example.com>
S: 250 OK
C: RCPT TO:<bob@example.org>
S: 250 OK
C: DATA
S: 354 End data with <CR><LF>.<CR><LF>
C: Subject: Test
C:
C: Hello, Bob!
C: .
S: 250 Message accepted
C: QUIT
S: 221 Bye
\end{lstlisting}

Здесь:
\begin{itemize}
    \item \texttt{EHLO} --- приветствие клиента;
    \item \texttt{MAIL FROM} --- адрес отправителя в SMTP-конверте;
    \item \texttt{RCPT TO} --- адрес получателя;
    \item \texttt{DATA} --- начало передачи содержимого письма;
    \item точка на отдельной строке завершает тело письма.
\end{itemize}

\subsection{POP3}

\textbf{POP3} --- Post Office Protocol version 3. Протокол предназначен для получения писем с почтового сервера.

Типичная модель POP3:
\begin{itemize}
    \item клиент подключается к серверу;
    \item проходит аутентификацию;
    \item загружает письма;
    \item при необходимости удаляет их с сервера;
    \item завершает сеанс.
\end{itemize}

Пример POP3-сессии:
\begin{lstlisting}
S: +OK POP3 server ready
C: USER student
S: +OK
C: PASS secret
S: +OK mailbox locked and ready
C: LIST
S: +OK 2 messages
S: 1 1200
S: 2 980
S: .
C: RETR 1
S: +OK message follows
S: ...
S: .
C: QUIT
S: +OK bye
\end{lstlisting}

\subsection{IMAP}

\textbf{IMAP} --- Internet Message Access Protocol. В отличие от POP3, IMAP ориентирован на хранение писем на сервере и синхронизацию состояния между несколькими устройствами.

IMAP поддерживает:
\begin{itemize}
    \item папки;
    \item флаги сообщений, например \texttt{Seen}, \texttt{Answered};
    \item выборочную загрузку частей письма;
    \item поиск по почтовому ящику;
    \item параллельную работу с нескольких устройств.
\end{itemize}

Пример IMAP-диалога:
\begin{lstlisting}
C: A001 LOGIN student secret
S: A001 OK LOGIN completed
C: A002 SELECT INBOX
S: * 3 EXISTS
S: A002 OK SELECT completed
C: A003 FETCH 1 BODY[HEADER]
S: * 1 FETCH ...
S: A003 OK FETCH completed
\end{lstlisting}

\subsection{Формат электронного письма}

Письмо состоит из:
\begin{enumerate}
    \item заголовков;
    \item пустой строки;
    \item тела сообщения.
\end{enumerate}

Пример:
\begin{lstlisting}
From: Alice <alice@example.com>
To: Bob <bob@example.org>
Subject: Meeting
Date: Wed, 6 May 2026 10:00:00 +0300
Content-Type: text/plain; charset=utf-8

Hello, Bob!
Let's meet tomorrow.
\end{lstlisting}

\subsection{MIME}

\textbf{MIME} --- Multipurpose Internet Mail Extensions. MIME расширяет возможности электронной почты и позволяет передавать:
\begin{itemize}
    \item текст в разных кодировках;
    \item HTML-сообщения;
    \item изображения;
    \item вложения;
    \item составные сообщения.
\end{itemize}

Пример MIME-заголовков:
\begin{lstlisting}
MIME-Version: 1.0
Content-Type: multipart/mixed; boundary="boundary123"
\end{lstlisting}

\subsection{Алгоритм доставки электронного письма}

Пусть пользователь Alice отправляет письмо пользователю Bob на адрес:
\[
\texttt{bob@example.org}
\]

Алгоритм:
\begin{enumerate}
    \item Alice составляет письмо в MUA.
    \item MUA передаёт письмо на SMTP-сервер отправителя.
    \item SMTP-сервер определяет домен получателя: \texttt{example.org}.
    \item Сервер отправителя выполняет DNS-запрос MX-записи для домена \texttt{example.org}.
    \item Получив адрес почтового сервера домена, сервер отправителя открывает SMTP-соединение с сервером получателя.
    \item Письмо передаётся по SMTP.
    \item Сервер получателя принимает письмо и передаёт его MDA.
    \item MDA помещает письмо в почтовый ящик Bob.
    \item Bob получает доступ к письму через IMAP, POP3 или веб-интерфейс.
\end{enumerate}

Схема:
\[
\text{MUA}_A \rightarrow \text{SMTP}_A \rightarrow \text{DNS/MX} \rightarrow \text{SMTP}_B \rightarrow \text{Mailbox}_B \rightarrow \text{MUA}_B
\]

\subsection{Безопасность электронной почты}

Основные проблемы:
\begin{itemize}
    \item перехват писем;
    \item подделка адреса отправителя;
    \item спам;
    \item фишинг;
    \item вредоносные вложения.
\end{itemize}

Средства защиты:
\begin{itemize}
    \item TLS для шифрования соединения;
    \item SMTP AUTH для аутентификации отправителя;
    \item SPF для указания разрешённых серверов отправки;
    \item DKIM для криптографической подписи писем доменом;
    \item DMARC для политики проверки SPF и DKIM;
    \item S/MIME или OpenPGP для сквозного шифрования и подписи сообщений.
\end{itemize}

\section{Телеконференции и группы новостей}

\subsection{Понятие телеконференции}

\textbf{Телеконференция} --- форма сетевого общения, при которой сообщения распространяются среди группы участников, объединённых общей темой.

Исторически важной системой телеконференций была \textbf{Usenet}.

\subsection{Usenet}

\textbf{Usenet} --- распределённая система обмена сообщениями в тематических группах новостей.

Сообщения Usenet организуются в группы, например:
\begin{lstlisting}
comp.lang.c
sci.math
rec.music.classical
\end{lstlisting}

Иерархии:
\begin{itemize}
    \item \texttt{comp} --- компьютерная тематика;
    \item \texttt{sci} --- наука;
    \item \texttt{rec} --- отдых и развлечения;
    \item \texttt{news} --- обсуждение самой системы новостей;
    \item \texttt{soc} --- социальные темы;
    \item \texttt{talk} --- дискуссионные темы.
\end{itemize}

\subsection{NNTP}

\textbf{NNTP} --- Network News Transfer Protocol, протокол передачи сетевых новостей.

NNTP используется:
\begin{itemize}
    \item для чтения сообщений групп новостей;
    \item для публикации новых сообщений;
    \item для обмена сообщениями между news-серверами.
\end{itemize}

Типичные команды NNTP:
\begin{center}
\begin{tabular}{|c|p{8cm}|}
\hline
Команда & Назначение \\
\hline
\texttt{GROUP} & выбор группы новостей \\
\hline
\texttt{ARTICLE} & получение статьи \\
\hline
\texttt{HEAD} & получение заголовков статьи \\
\hline
\texttt{BODY} & получение тела статьи \\
\hline
\texttt{POST} & публикация статьи \\
\hline
\texttt{QUIT} & завершение соединения \\
\hline
\end{tabular}
\end{center}

Пример NNTP-сессии:
\begin{lstlisting}
S: 200 news.example.org NNTP server ready
C: GROUP comp.lang.c
S: 211 123 1 123 comp.lang.c
C: ARTICLE 10
S: 220 10 <msgid@example.org> article follows
S: From: user@example.org
S: Subject: Question about pointers
S:
S: ...
S: .
C: QUIT
S: 205 closing connection
\end{lstlisting}

\subsection{Сравнение электронной почты и телеконференций}

\begin{center}
\begin{tabular}{|p{4cm}|p{5cm}|p{5cm}|}
\hline
Критерий & Электронная почта & Телеконференции/Usenet \\
\hline
Модель общения & адресная: от отправителя к получателям & тематическая: сообщения публикуются в группе \\
\hline
Хранение & в почтовых ящиках пользователей & на news-серверах \\
\hline
Получатели & явно указаны в письме & все подписчики или читатели группы \\
\hline
Протоколы & SMTP, POP3, IMAP & NNTP \\
\hline
Основное назначение & личная и деловая переписка & публичные обсуждения по темам \\
\hline
\end{tabular}
\end{center}

\section{Протокол FTP}

\subsection{Назначение FTP}

\textbf{FTP} --- File Transfer Protocol, протокол передачи файлов. Он предназначен для обмена файлами между клиентом и сервером.

FTP позволяет:
\begin{itemize}
    \item просматривать каталоги на сервере;
    \item загружать файлы с сервера;
    \item отправлять файлы на сервер;
    \item удалять и переименовывать файлы;
    \item создавать каталоги.
\end{itemize}

\subsection{Архитектура FTP}

FTP использует два соединения:
\begin{enumerate}
    \item \textbf{управляющее соединение} --- для команд и ответов;
    \item \textbf{соединение данных} --- для передачи содержимого файлов и списков каталогов.
\end{enumerate}

Обычно управляющее соединение устанавливается на порт:
\[
21
\]

Соединение данных в классическом активном режиме использует порт:
\[
20
\]

\subsection{Активный и пассивный режимы FTP}

\subsubsection{Активный режим}

В активном режиме:
\begin{enumerate}
    \item клиент подключается к серверу на порт 21;
    \item клиент сообщает серверу свой адрес и порт командой \texttt{PORT};
    \item сервер сам открывает соединение данных к клиенту.
\end{enumerate}

Схема:
\[
\text{Client} \xrightarrow{\text{control}} \text{Server:21}
\]
\[
\text{Server:20} \xrightarrow{\text{data}} \text{Client}
\]

Недостаток: активный режим плохо работает через NAT и межсетевые экраны, поскольку сервер должен открыть входящее соединение к клиенту.

\subsubsection{Пассивный режим}

В пассивном режиме:
\begin{enumerate}
    \item клиент подключается к серверу на порт 21;
    \item клиент отправляет команду \texttt{PASV};
    \item сервер сообщает клиенту IP-адрес и порт для соединения данных;
    \item клиент сам открывает соединение данных к серверу.
\end{enumerate}

Схема:
\[
\text{Client} \xrightarrow{\text{control}} \text{Server:21}
\]
\[
\text{Client} \xrightarrow{\text{data}} \text{Server:passive port}
\]

Пассивный режим обычно удобнее при работе через NAT.

\subsection{Основные FTP-команды}

\begin{center}
\begin{tabular}{|c|p{9cm}|}
\hline
Команда & Назначение \\
\hline
\texttt{USER} & передача имени пользователя \\
\hline
\texttt{PASS} & передача пароля \\
\hline
\texttt{PWD} & вывод текущего каталога \\
\hline
\texttt{CWD} & смена каталога \\
\hline
\texttt{LIST} & получение списка файлов \\
\hline
\texttt{RETR} & загрузка файла с сервера \\
\hline
\texttt{STOR} & отправка файла на сервер \\
\hline
\texttt{DELE} & удаление файла \\
\hline
\texttt{MKD} & создание каталога \\
\hline
\texttt{RMD} & удаление каталога \\
\hline
\texttt{QUIT} & завершение сеанса \\
\hline
\end{tabular}
\end{center}

\subsection{Коды ответов FTP}

FTP-сервер возвращает числовые коды ответов.

\begin{center}
\begin{tabular}{|c|p{9cm}|}
\hline
Код & Значение \\
\hline
\texttt{1xx} & предварительный положительный ответ \\
\hline
\texttt{2xx} & успешное выполнение команды \\
\hline
\texttt{3xx} & команда принята, но требуется дополнительная информация \\
\hline
\texttt{4xx} & временная ошибка \\
\hline
\texttt{5xx} & постоянная ошибка \\
\hline
\end{tabular}
\end{center}

Примеры:
\begin{lstlisting}
220 Service ready
331 User name okay, need password
230 User logged in
550 File unavailable
\end{lstlisting}

\subsection{Пример FTP-сеанса}

\begin{lstlisting}
S: 220 FTP server ready
C: USER student
S: 331 Password required
C: PASS secret
S: 230 User logged in
C: PWD
S: 257 "/home/student" is current directory
C: PASV
S: 227 Entering Passive Mode (192,0,2,1,195,80)
C: LIST
S: 150 Opening data connection
S: 226 Transfer complete
C: RETR report.txt
S: 150 Opening data connection
S: 226 Transfer complete
C: QUIT
S: 221 Goodbye
\end{lstlisting}

\subsection{Алгоритм загрузки файла по FTP в пассивном режиме}

Пусть клиент хочет скачать файл \texttt{data.txt}.

\begin{enumerate}
    \item Клиент открывает управляющее TCP-соединение с сервером на порт 21.
    \item Сервер отправляет приветствие.
    \item Клиент отправляет \texttt{USER}.
    \item Клиент отправляет \texttt{PASS}.
    \item Сервер подтверждает вход.
    \item Клиент отправляет \texttt{PASV}.
    \item Сервер сообщает порт для соединения данных.
    \item Клиент открывает второе TCP-соединение к указанному порту.
    \item Клиент отправляет команду \texttt{RETR data.txt}.
    \item Сервер передаёт содержимое файла по соединению данных.
    \item После передачи соединение данных закрывается.
    \item Управляющее соединение может использоваться для следующих команд или закрывается командой \texttt{QUIT}.
\end{enumerate}

\subsection{Недостатки FTP}

Классический FTP имеет существенные недостатки:
\begin{itemize}
    \item логины и пароли передаются в открытом виде;
    \item данные передаются без шифрования;
    \item отдельное соединение данных осложняет работу через NAT и firewall;
    \item отсутствует встроенная криптографическая проверка подлинности сервера.
\end{itemize}

Более безопасные альтернативы:
\begin{itemize}
    \item \textbf{FTPS} --- FTP поверх TLS;
    \item \textbf{SFTP} --- протокол передачи файлов поверх SSH, не являющийся FTP;
    \item \textbf{HTTPS} --- часто используется для скачивания файлов;
    \item \textbf{rsync over SSH} --- эффективная синхронизация файлов.
\end{itemize}

\section{Протокол HTTP}

\subsection{Назначение HTTP}

\textbf{HTTP} --- Hypertext Transfer Protocol, прикладной протокол передачи гипертекстовых документов и других ресурсов.

HTTP используется для:
\begin{itemize}
    \item загрузки HTML-страниц;
    \item получения изображений, CSS, JavaScript и других ресурсов;
    \item взаимодействия веб-клиентов с веб-приложениями;
    \item построения REST API;
    \item передачи файлов.
\end{itemize}

HTTP является протоколом прикладного уровня и обычно работает поверх TCP. HTTPS --- это HTTP, работающий поверх защищённого канала TLS.

\subsection{Основные свойства HTTP}

\begin{enumerate}
    \item \textbf{Модель запрос--ответ.} Клиент отправляет запрос, сервер возвращает ответ.
    \item \textbf{Отсутствие состояния на уровне протокола.} Каждый запрос самодостаточен; сервер не обязан помнить предыдущие запросы.
    \item \textbf{Расширяемость.} Заголовки позволяют добавлять новые возможности.
    \item \textbf{Универсальность представлений.} Через HTTP можно передавать HTML, JSON, XML, изображения, видео, архивы и др.
\end{enumerate}

\subsection{HTTP-запрос}

HTTP-запрос состоит из:
\begin{enumerate}
    \item стартовой строки;
    \item заголовков;
    \item пустой строки;
    \item необязательного тела.
\end{enumerate}

Пример:
\begin{lstlisting}
GET /index.html HTTP/1.1
Host: example.org
User-Agent: SimpleBrowser/1.0
Accept: text/html

\end{lstlisting}

Стартовая строка:
\[
\texttt{GET /index.html HTTP/1.1}
\]

содержит:
\begin{itemize}
    \item метод \texttt{GET};
    \item путь \texttt{/index.html};
    \item версию протокола \texttt{HTTP/1.1}.
\end{itemize}

\subsection{HTTP-ответ}

HTTP-ответ состоит из:
\begin{enumerate}
    \item строки состояния;
    \item заголовков;
    \item пустой строки;
    \item тела ответа.
\end{enumerate}

Пример:
\begin{lstlisting}
HTTP/1.1 200 OK
Content-Type: text/html; charset=utf-8
Content-Length: 72

<!doctype html>
<html>
  <body>Hello, world!</body>
</html>
\end{lstlisting}

\subsection{Методы HTTP}

\begin{center}
\begin{tabular}{|c|p{9cm}|}
\hline
Метод & Назначение \\
\hline
\texttt{GET} & получение ресурса \\
\hline
\texttt{HEAD} & получение только заголовков ответа \\
\hline
\texttt{POST} & отправка данных на сервер, часто создание подчинённого ресурса или выполнение действия \\
\hline
\texttt{PUT} & замена или создание ресурса по заданному URI \\
\hline
\texttt{PATCH} & частичное изменение ресурса \\
\hline
\texttt{DELETE} & удаление ресурса \\
\hline
\texttt{OPTIONS} & получение информации о поддерживаемых возможностях \\
\hline
\texttt{CONNECT} & установление туннеля, например для HTTPS через прокси \\
\hline
\texttt{TRACE} & диагностический метод, обычно отключается из соображений безопасности \\
\hline
\end{tabular}
\end{center}

\subsection{Безопасные и идемпотентные методы}

\textbf{Безопасный метод} --- метод, который не должен изменять состояние сервера с точки зрения семантики запроса. Примеры: \texttt{GET}, \texttt{HEAD}, \texttt{OPTIONS}.

\textbf{Идемпотентный метод} --- метод, повторное выполнение которого имеет тот же эффект, что и однократное выполнение.

Идемпотентные методы:
\[
\texttt{GET}, \texttt{HEAD}, \texttt{PUT}, \texttt{DELETE}, \texttt{OPTIONS}
\]

Метод \texttt{POST} в общем случае не является идемпотентным.

Пример:
\begin{itemize}
    \item \texttt{PUT /users/10} с одним и тем же телом несколько раз приводит к одному и тому же состоянию ресурса;
    \item \texttt{POST /orders} несколько раз может создать несколько разных заказов.
\end{itemize}

\subsection{Коды состояния HTTP}

\begin{center}
\begin{tabular}{|c|p{9cm}|}
\hline
Класс & Значение \\
\hline
\texttt{1xx} & информационные ответы \\
\hline
\texttt{2xx} & успешная обработка запроса \\
\hline
\texttt{3xx} & перенаправление \\
\hline
\texttt{4xx} & ошибка клиента \\
\hline
\texttt{5xx} & ошибка сервера \\
\hline
\end{tabular}
\end{center}

Часто используемые коды:
\begin{center}
\begin{tabular}{|c|p{9cm}|}
\hline
Код & Значение \\
\hline
\texttt{200 OK} & запрос успешно выполнен \\
\hline
\texttt{201 Created} & ресурс создан \\
\hline
\texttt{301 Moved Permanently} & ресурс перемещён постоянно \\
\hline
\texttt{302 Found} & временное перенаправление \\
\hline
\texttt{304 Not Modified} & ресурс не изменился, можно использовать кэш \\
\hline
\texttt{400 Bad Request} & некорректный запрос \\
\hline
\texttt{401 Unauthorized} & требуется аутентификация \\
\hline
\texttt{403 Forbidden} & доступ запрещён \\
\hline
\texttt{404 Not Found} & ресурс не найден \\
\hline
\texttt{500 Internal Server Error} & внутренняя ошибка сервера \\
\hline
\texttt{502 Bad Gateway} & ошибка шлюза \\
\hline
\texttt{503 Service Unavailable} & сервис временно недоступен \\
\hline
\end{tabular}
\end{center}

\subsection{Заголовки HTTP}

HTTP-заголовки передают метаданные запроса или ответа.

Примеры заголовков запроса:
\begin{lstlisting}
Host: example.org
User-Agent: Mozilla/5.0
Accept: text/html
Accept-Language: ru
Cookie: sessionid=abc123
\end{lstlisting}

Примеры заголовков ответа:
\begin{lstlisting}
Content-Type: text/html; charset=utf-8
Content-Length: 1024
Set-Cookie: sessionid=abc123; HttpOnly; Secure
Cache-Control: max-age=3600
Location: https://example.org/new-page
\end{lstlisting}

\subsection{Cookies}

\textbf{Cookie} --- небольшой фрагмент данных, который сервер передаёт браузеру, а браузер затем отправляет обратно этому серверу.

Cookies используются для:
\begin{itemize}
    \item хранения идентификатора сессии;
    \item сохранения пользовательских настроек;
    \item аналитики;
    \item авторизации.
\end{itemize}

Пример установки cookie:
\begin{lstlisting}
Set-Cookie: sessionid=abc123; HttpOnly; Secure; SameSite=Lax
\end{lstlisting}

Важные атрибуты:
\begin{itemize}
    \item \texttt{HttpOnly} --- запрещает доступ к cookie из JavaScript;
    \item \texttt{Secure} --- cookie передаётся только по HTTPS;
    \item \texttt{SameSite} --- ограничивает отправку cookie при межсайтовых запросах.
\end{itemize}

\subsection{Кэширование HTTP}

HTTP поддерживает кэширование для уменьшения задержек и нагрузки на сервер.

Основные механизмы:
\begin{itemize}
    \item \texttt{Cache-Control};
    \item \texttt{ETag};
    \item \texttt{Last-Modified};
    \item условные запросы с \texttt{If-None-Match} и \texttt{If-Modified-Since}.
\end{itemize}

Пример:
\begin{lstlisting}
GET /style.css HTTP/1.1
Host: example.org
If-None-Match: "abc123"
\end{lstlisting}

Ответ:
\begin{lstlisting}
HTTP/1.1 304 Not Modified
ETag: "abc123"
\end{lstlisting}

В этом случае тело ресурса не передаётся, браузер использует локальную копию.

\subsection{Алгоритм загрузки веб-страницы браузером}

Пусть пользователь вводит:
\begin{lstlisting}
https://example.org/index.html
\end{lstlisting}

Алгоритм:
\begin{enumerate}
    \item Браузер разбирает URL.
    \item Проверяет кэш, HSTS-политику и другие локальные данные.
    \item Выполняет DNS-разрешение имени \texttt{example.org}.
    \item Устанавливает TCP-соединение с сервером.
    \item Для HTTPS выполняет TLS-рукопожатие.
    \item Формирует HTTP-запрос.
    \item Отправляет запрос серверу.
    \item Сервер обрабатывает запрос.
    \item Сервер возвращает HTTP-ответ.
    \item Браузер анализирует заголовки ответа.
    \item Браузер разбирает HTML.
    \item При обнаружении внешних ресурсов, например CSS, JS, изображений, браузер отправляет дополнительные запросы.
    \item Строятся DOM, CSSOM и дерево рендеринга.
    \item Выполняется раскладка элементов на странице.
    \item Выполняется отрисовка.
\end{enumerate}

\subsection{HTTPS и TLS}

\textbf{HTTPS} --- это использование HTTP поверх TLS.

TLS обеспечивает:
\begin{itemize}
    \item конфиденциальность передаваемых данных;
    \item целостность;
    \item аутентификацию сервера с помощью сертификата;
    \item при необходимости аутентификацию клиента.
\end{itemize}

Упрощённые этапы TLS-соединения:
\begin{enumerate}
    \item Клиент отправляет \texttt{ClientHello}: поддерживаемые версии TLS, наборы шифров, случайные данные.
    \item Сервер отправляет \texttt{ServerHello}: выбранные параметры.
    \item Сервер передаёт сертификат.
    \item Клиент проверяет сертификат.
    \item Стороны согласуют общий секрет.
    \item Из общего секрета выводятся симметричные ключи.
    \item Дальнейший HTTP-трафик передаётся в зашифрованном виде.
\end{enumerate}

\section{Язык HTML}

\subsection{Понятие HTML}

\textbf{HTML} --- HyperText Markup Language, язык гипертекстовой разметки. HTML используется для описания структуры веб-документов.

HTML не является языком программирования, поскольку не задаёт алгоритмы вычислений. Он описывает смысловую и структурную организацию документа.

\subsection{HTML-документ}

Минимальный HTML-документ:
\begin{lstlisting}[language=HTML]
<!doctype html>
<html lang="ru">
<head>
    <meta charset="utf-8">
    <title>Минимальная страница</title>
</head>
<body>
    <h1>Привет, мир!</h1>
    <p>Это пример HTML-документа.</p>
</body>
</html>
\end{lstlisting}

Назначение элементов:
\begin{itemize}
    \item \texttt{<!doctype html>} --- объявление типа документа;
    \item \texttt{html} --- корневой элемент;
    \item \texttt{head} --- служебная информация;
    \item \texttt{meta charset} --- кодировка;
    \item \texttt{title} --- заголовок вкладки браузера;
    \item \texttt{body} --- видимое содержимое страницы.
\end{itemize}

\subsection{Элементы и атрибуты}

HTML-документ состоит из элементов.

Общий вид элемента:
\begin{lstlisting}[language=HTML]
<tag attribute="value">content</tag>
\end{lstlisting}

Пример:
\begin{lstlisting}[language=HTML]
<a href="https://example.org">Ссылка</a>
\end{lstlisting}

Здесь:
\begin{itemize}
    \item \texttt{a} --- элемент ссылки;
    \item \texttt{href} --- атрибут;
    \item \texttt{https://example.org} --- значение атрибута;
    \item \texttt{Ссылка} --- содержимое элемента.
\end{itemize}

\subsection{Основные группы HTML-элементов}

\subsubsection{Структурные элементы}

\begin{lstlisting}[language=HTML]
<header>Шапка сайта</header>
<nav>Навигация</nav>
<main>Основное содержимое</main>
<section>Раздел</section>
<article>Самостоятельная статья</article>
<aside>Боковая информация</aside>
<footer>Подвал</footer>
\end{lstlisting}

\subsubsection{Заголовки и текст}

\begin{lstlisting}[language=HTML]
<h1>Главный заголовок</h1>
<h2>Подзаголовок</h2>
<p>Абзац текста.</p>
<strong>Важный текст</strong>
<em>Смысловое выделение</em>
\end{lstlisting}

\subsubsection{Списки}

Нумерованный список:
\begin{lstlisting}[language=HTML]
<ol>
    <li>Первый пункт</li>
    <li>Второй пункт</li>
</ol>
\end{lstlisting}

Маркированный список:
\begin{lstlisting}[language=HTML]
<ul>
    <li>Элемент</li>
    <li>Ещё элемент</li>
</ul>
\end{lstlisting}

\subsubsection{Таблицы}

\begin{lstlisting}[language=HTML]
<table>
    <tr>
        <th>Имя</th>
        <th>Оценка</th>
    </tr>
    <tr>
        <td>Анна</td>
        <td>5</td>
    </tr>
</table>
\end{lstlisting}

\subsubsection{Формы}

\begin{lstlisting}[language=HTML]
<form action="/login" method="post">
    <label>
        Логин:
        <input type="text" name="username">
    </label>
    <label>
        Пароль:
        <input type="password" name="password">
    </label>
    <button type="submit">Войти</button>
</form>
\end{lstlisting}

Элемент \texttt{form} задаёт форму, а атрибуты:
\begin{itemize}
    \item \texttt{action} --- адрес отправки данных;
    \item \texttt{method} --- HTTP-метод отправки.
\end{itemize}

\subsection{Гипертекстовые ссылки}

Гипертекстовая ссылка создаётся элементом \texttt{a}:
\begin{lstlisting}[language=HTML]
<a href="page.html">Перейти на страницу</a>
\end{lstlisting}

Абсолютная ссылка:
\begin{lstlisting}[language=HTML]
<a href="https://example.org/docs/page.html">Документация</a>
\end{lstlisting}

Относительная ссылка:
\begin{lstlisting}[language=HTML]
<a href="../index.html">На главную</a>
\end{lstlisting}

Ссылка на фрагмент страницы:
\begin{lstlisting}[language=HTML]
<a href="#chapter1">К главе 1</a>

<h2 id="chapter1">Глава 1</h2>
\end{lstlisting}

\subsection{HTML и DOM}

\textbf{DOM} --- Document Object Model, объектная модель документа. Браузер разбирает HTML-код и строит дерево DOM.

Пример HTML:
\begin{lstlisting}[language=HTML]
<html>
  <body>
    <h1>Title</h1>
    <p>Hello</p>
  </body>
</html>
\end{lstlisting}

Соответствующее дерево:
\[
\texttt{html}
\rightarrow
\texttt{body}
\rightarrow
\{\texttt{h1}, \texttt{p}\}
\]

DOM позволяет JavaScript изменять документ во время работы страницы:
\begin{lstlisting}[language=JavaScript]
document.querySelector("h1").textContent = "Новый заголовок";
\end{lstlisting}

\subsection{Семантическая разметка}

\textbf{Семантическая разметка} --- использование HTML-элементов в соответствии с их смыслом.

Например, правильно:
\begin{lstlisting}[language=HTML]
<nav>
    <a href="/">Главная</a>
    <a href="/about">О сайте</a>
</nav>
\end{lstlisting}

Менее правильно:
\begin{lstlisting}[language=HTML]
<div>
    <span>Главная</span>
    <span>О сайте</span>
</div>
\end{lstlisting}

Преимущества семантической разметки:
\begin{itemize}
    \item улучшение доступности;
    \item лучшая индексация поисковыми системами;
    \item понятная структура документа;
    \item облегчение сопровождения кода.
\end{itemize}

\section{Разработка Web-страниц}

\subsection{Основные технологии Web}

Современная веб-страница обычно строится на трёх базовых технологиях:
\begin{enumerate}
    \item \textbf{HTML} --- структура документа;
    \item \textbf{CSS} --- внешний вид;
    \item \textbf{JavaScript} --- поведение и интерактивность.
\end{enumerate}

\subsection{CSS}

\textbf{CSS} --- Cascading Style Sheets, каскадные таблицы стилей.

Пример:
\begin{lstlisting}[language=HTML]
<!doctype html>
<html lang="ru">
<head>
    <meta charset="utf-8">
    <title>CSS example</title>
    <style>
        body {
            font-family: sans-serif;
            background: #f5f5f5;
        }

        h1 {
            color: navy;
        }
    </style>
</head>
<body>
    <h1>Заголовок</h1>
    <p>Текст страницы.</p>
</body>
</html>
\end{lstlisting}

\subsection{Подключение внешнего CSS}

\begin{lstlisting}[language=HTML]
<link rel="stylesheet" href="style.css">
\end{lstlisting}

Файл \texttt{style.css}:
\begin{lstlisting}[language=CSS]
body {
    font-family: Arial, sans-serif;
}

main {
    max-width: 800px;
    margin: 0 auto;
}
\end{lstlisting}

\subsection{JavaScript}

\textbf{JavaScript} --- язык программирования, используемый в браузере для создания интерактивных веб-страниц.

Пример:
\begin{lstlisting}[language=HTML]
<button id="btn">Нажми меня</button>
<p id="result"></p>

<script>
    const button = document.getElementById("btn");
    const result = document.getElementById("result");

    button.addEventListener("click", function () {
        result.textContent = "Кнопка нажата";
    });
</script>
\end{lstlisting}

\subsection{Статические и динамические страницы}

\textbf{Статическая веб-страница} --- файл, который сервер отдаёт клиенту без генерации на стороне сервера.

Пример:
\begin{lstlisting}
index.html
style.css
script.js
\end{lstlisting}

\textbf{Динамическая веб-страница} --- страница, содержимое которой формируется программой, часто с использованием базы данных.

Примеры серверных технологий:
\begin{itemize}
    \item PHP;
    \item Python с Django или Flask;
    \item Java с Spring;
    \item JavaScript/TypeScript с Node.js;
    \item Ruby on Rails;
    \item Go;
    \item ASP.NET.
\end{itemize}

\subsection{Минимальный пример сайта}

Структура проекта:
\begin{lstlisting}
site/
  index.html
  style.css
  script.js
\end{lstlisting}

Файл \texttt{index.html}:
\begin{lstlisting}[language=HTML]
<!doctype html>
<html lang="ru">
<head>
    <meta charset="utf-8">
    <title>Мини-сайт</title>
    <link rel="stylesheet" href="style.css">
</head>
<body>
    <main>
        <h1>Мини-сайт</h1>
        <p id="message">Добро пожаловать!</p>
        <button id="change">Изменить текст</button>
    </main>
    <script src="script.js"></script>
</body>
</html>
\end{lstlisting}

Файл \texttt{style.css}:
\begin{lstlisting}[language=CSS]
body {
    font-family: sans-serif;
    background: #eeeeee;
}

main {
    max-width: 700px;
    margin: 40px auto;
    padding: 20px;
    background: white;
    border-radius: 8px;
}
\end{lstlisting}

Файл \texttt{script.js}:
\begin{lstlisting}[language=JavaScript]
document.getElementById("change").addEventListener("click", () => {
    document.getElementById("message").textContent = "Текст изменён!";
});
\end{lstlisting}

\subsection{Этапы разработки Web-страницы}

\begin{enumerate}
    \item Определить цель страницы.
    \item Спроектировать структуру содержимого.
    \item Написать HTML-разметку.
    \item Добавить CSS-стили.
    \item Добавить JavaScript при необходимости.
    \item Проверить отображение в разных браузерах.
    \item Проверить адаптивность на разных размерах экрана.
    \item Проверить доступность.
    \item Оптимизировать изображения и ресурсы.
    \item Разместить сайт на веб-сервере.
\end{enumerate}

\subsection{Адаптивная вёрстка}

\textbf{Адаптивная вёрстка} --- подход, при котором страница корректно отображается на устройствах с разной шириной экрана.

Пример CSS media query:
\begin{lstlisting}[language=CSS]
.container {
    display: flex;
    gap: 20px;
}

@media (max-width: 600px) {
    .container {
        flex-direction: column;
    }
}
\end{lstlisting}

\subsection{Доступность Web-страниц}

\textbf{Доступность} означает, что сайтом могут пользоваться люди с различными ограничениями возможностей.

Основные рекомендации:
\begin{itemize}
    \item использовать семантические HTML-элементы;
    \item указывать альтернативный текст для изображений;
    \item обеспечивать достаточный контраст;
    \item делать сайт доступным с клавиатуры;
    \item правильно связывать \texttt{label} и \texttt{input};
    \item не полагаться только на цвет для передачи смысла.
\end{itemize}

Пример:
\begin{lstlisting}[language=HTML]
<label for="email">Email</label>
<input id="email" name="email" type="email">
\end{lstlisting}

\section{WWW и WWW-серверы}

\subsection{World Wide Web}

\textbf{World Wide Web} --- глобальная распределённая гипертекстовая информационная система, работающая поверх Internet.

Не следует отождествлять Internet и Web:
\begin{itemize}
    \item \textbf{Internet} --- глобальная сеть сетей, инфраструктура передачи данных;
    \item \textbf{WWW} --- один из сервисов Internet, основанный на HTTP, URL и HTML.
\end{itemize}

Основные компоненты WWW:
\begin{enumerate}
    \item URL/URI --- идентификация ресурсов;
    \item HTTP/HTTPS --- протокол доступа;
    \item HTML --- формат гипертекстовых документов;
    \item браузеры --- клиенты Web;
    \item web-серверы --- поставщики ресурсов.
\end{enumerate}

\subsection{WWW-сервер}

\textbf{WWW-сервер} или \textbf{Web-сервер} --- программа, принимающая HTTP-запросы и возвращающая HTTP-ответы.

Популярные Web-серверы:
\begin{itemize}
    \item Apache HTTP Server;
    \item Nginx;
    \item Microsoft IIS;
    \item Caddy;
    \item Lighttpd.
\end{itemize}

\subsection{Основные функции Web-сервера}

Web-сервер выполняет:
\begin{itemize}
    \item приём TCP-соединений;
    \item приём и разбор HTTP-запросов;
    \item определение нужного ресурса;
    \item выдачу статических файлов;
    \item передачу запросов приложению через FastCGI, WSGI, reverse proxy и др.;
    \item формирование HTTP-ответов;
    \item ведение журналов доступа и ошибок;
    \item сжатие данных;
    \item TLS-терминацию;
    \item балансировку нагрузки;
    \item кэширование;
    \item ограничение доступа.
\end{itemize}

\subsection{Статический и динамический Web-сервер}

\textbf{Статическая выдача}:
\[
\text{HTTP-запрос} \rightarrow \text{файл на диске} \rightarrow \text{HTTP-ответ}
\]

\textbf{Динамическая выдача}:
\[
\text{HTTP-запрос} \rightarrow \text{серверное приложение} \rightarrow \text{БД} \rightarrow \text{HTML/JSON-ответ}
\]

\subsection{CGI}

\textbf{CGI} --- Common Gateway Interface, ранний стандарт взаимодействия Web-сервера с внешними программами.

Алгоритм работы CGI:
\begin{enumerate}
    \item Web-сервер получает HTTP-запрос.
    \item Определяет, что запрос должен быть обработан CGI-программой.
    \item Запускает внешний процесс.
    \item Передаёт ему параметры запроса через переменные окружения и стандартный ввод.
    \item CGI-программа формирует ответ.
    \item Web-сервер передаёт ответ клиенту.
\end{enumerate}

Недостаток CGI: запуск отдельного процесса на каждый запрос создаёт большие накладные расходы.

\subsection{Reverse proxy}

\textbf{Обратный прокси-сервер} принимает запросы от клиентов и перенаправляет их внутренним серверам.

Схема:
\[
\text{Client} \rightarrow \text{Nginx} \rightarrow \text{Application Server}
\]

Задачи reverse proxy:
\begin{itemize}
    \item TLS-терминация;
    \item балансировка нагрузки;
    \item кэширование;
    \item защита внутренних серверов;
    \item маршрутизация запросов;
    \item сжатие ответов.
\end{itemize}

\subsection{Минимальный учебный HTTP-сервер}

Ниже приведён минимальный пример HTTP-сервера на Python. Он принимает соединение, читает запрос и возвращает HTML-страницу.

\begin{lstlisting}[language=Python]
import socket

HOST = "127.0.0.1"
PORT = 8080

server = socket.socket(socket.AF_INET, socket.SOCK_STREAM)
server.bind((HOST, PORT))
server.listen(1)

print(f"Server started at http://{HOST}:{PORT}")

while True:
    conn, addr = server.accept()
    request = conn.recv(4096).decode("utf-8", errors="ignore")
    print(request)

    body = """<!doctype html>
<html>
<body>
<h1>Hello from simple server</h1>
</body>
</html>"""

    response = (
        "HTTP/1.1 200 OK\r\n"
        "Content-Type: text/html; charset=utf-8\r\n"
        f"Content-Length: {len(body.encode('utf-8'))}\r\n"
        "\r\n"
        f"{body}"
    )

    conn.sendall(response.encode("utf-8"))
    conn.close()
\end{lstlisting}

Этот пример иллюстрирует базовый принцип HTTP, но не является полноценным сервером: он не обрабатывает параллельные соединения, ошибки, TLS, разные методы и MIME-типы.

\section{Сравнение FTP и HTTP}

\begin{center}
\begin{tabular}{|p{4cm}|p{5cm}|p{5cm}|}
\hline
Критерий & FTP & HTTP \\
\hline
Основное назначение & Передача файлов & Передача гипертекстовых и других ресурсов \\
\hline
Модель соединений & Управляющее соединение и соединение данных & Обычно одно логическое соединение для запросов/ответов \\
\hline
Состояние сеанса & Более выраженное состояние: вход, текущий каталог, режим передачи & Протокол сам по себе stateless \\
\hline
Типичные порты & 21, 20 & 80, 443 \\
\hline
Работа через NAT & Может быть сложной, особенно в активном режиме & Обычно проще \\
\hline
Защита & В классическом FTP отсутствует & HTTPS широко применяет TLS \\
\hline
Использование в Web & Ограниченное, историческое & Основной протокол Web \\
\hline
\end{tabular}
\end{center}

\section{Безопасность удалённого доступа и Web}

\subsection{Основные угрозы}

\begin{itemize}
    \item перехват трафика;
    \item подбор паролей;
    \item атаки ``человек посередине'';
    \item эксплуатация уязвимостей серверов;
    \item фишинг;
    \item межсайтовый скриптинг;
    \item SQL-инъекции;
    \item утечка сессионных cookie;
    \item неправильная настройка прав доступа.
\end{itemize}

\subsection{Общие меры защиты}

\begin{enumerate}
    \item Использовать защищённые протоколы: SSH вместо Telnet, HTTPS вместо HTTP для конфиденциальных данных.
    \item Применять сильную аутентификацию.
    \item Использовать многофакторную аутентификацию.
    \item Обновлять серверное ПО.
    \item Ограничивать доступ по IP или через VPN.
    \item Настраивать firewall.
    \item Хранить пароли в виде криптографических хэшей с солью.
    \item Проверять и фильтровать пользовательский ввод.
    \item Использовать принцип минимальных привилегий.
    \item Вести журналы событий и анализировать их.
\end{enumerate}

\subsection{Web-угрозы}

\subsubsection{XSS}

\textbf{XSS} --- Cross-Site Scripting, внедрение вредоносного JavaScript-кода в страницу.

Пример опасной ситуации:
\begin{lstlisting}[language=HTML]
<p>Здравствуйте, <script>alert('XSS')</script></p>
\end{lstlisting}

Защита:
\begin{itemize}
    \item экранирование HTML;
    \item Content Security Policy;
    \item запрет вставки непроверенного HTML;
    \item использование фреймворков с автоматическим escaping.
\end{itemize}

\subsubsection{SQL-инъекция}

SQL-инъекция возникает при включении пользовательского ввода непосредственно в SQL-запрос.

Плохой пример:
\begin{lstlisting}[language=Python]
query = "SELECT * FROM users WHERE name = '" + username + "'"
\end{lstlisting}

Правильный подход --- параметризованные запросы:
\begin{lstlisting}[language=Python]
cursor.execute(
    "SELECT * FROM users WHERE name = ?",
    (username,)
)
\end{lstlisting}

\subsubsection{CSRF}

\textbf{CSRF} --- Cross-Site Request Forgery, подделка межсайтового запроса. Суть атаки: браузер авторизованного пользователя автоматически отправляет запрос на доверенный сайт по инициативе вредоносной страницы.

Защита:
\begin{itemize}
    \item CSRF-токены;
    \item атрибут cookie \texttt{SameSite};
    \item проверка заголовков \texttt{Origin} и \texttt{Referer};
    \item требование повторного подтверждения важных операций.
\end{itemize}

\section{Итоговая связка технологий}

Работа типичной веб-системы включает несколько уровней:

\[
\text{Пользователь}
\rightarrow
\text{Браузер}
\rightarrow
\text{DNS}
\rightarrow
\text{TCP/TLS}
\rightarrow
\text{HTTP}
\rightarrow
\text{Web-сервер}
\rightarrow
\text{Приложение}
\rightarrow
\text{База данных}
\]

HTML, CSS и JavaScript формируют клиентскую часть. HTTP/HTTPS обеспечивает обмен данными между клиентом и сервером. Web-сервер принимает запросы, отдаёт статические ресурсы или передаёт запросы приложению. Электронная почта и телеконференции являются другими важными прикладными сервисами Internet, использующими собственные протоколы: SMTP, POP3, IMAP и NNTP. FTP исторически применялся для передачи файлов, но в современных системах часто заменяется более безопасными механизмами: SFTP, FTPS и HTTPS.

\end{document}
